\section{Prototipo funcional}

\subsection{Arquitectura verificada}

La evidencia operativa combina dos piezas. La primera es la demo reproducible
de esta etapa, que ejecuta la selección final sobre una imagen real. La segunda
es la aplicación móvil Doctor Maíz inspeccionada en el commit
\nolinkurl{dd048d37dcabd08c53c704d702825e7d672f0f37}. La aplicación usa React
Native 0.86, Expo 57, TypeScript, WatermelonDB y
\texttt{react-native-fast-tflite}. El modelo embarcado es
EfficientNet-Lite0 TFLite con pesos int8, entrada y salida float32 y tamaño de
3\,743\,824 bytes. Su SHA-256 coincide con el manifiesto:
\par\smallskip\noindent
\texttt{3a0623cd985a23b954e424e605a2e00c}\allowbreak
\texttt{91f8d592e376dfe3149e5820485bc2b3}.

\begin{figure}[H]
  \centering
  \fbox{\begin{minipage}{0.88\textwidth}
    \centering
    \textbf{Cámara o galería}
    $\longrightarrow$ orientación EXIF, RGB y $224\times224$
    $\longrightarrow$ \textbf{TFLite local}\\[0.3em]
    $\longrightarrow$ logits + características OOD
    $\longrightarrow$ Top-1, confianza y rechazo
    $\longrightarrow$ \textbf{resultado e historial offline}
  \end{minipage}}
  \caption{Flujo real implementado en la aplicación móvil. La API participa en
  autenticación, sincronización y actualización; la inferencia ocurre en el
  dispositivo.}
  \label{fig:prototype-flow}
\end{figure}

El motor materializa el modelo como archivo local, comprueba que el tensor de
entrada tenga $1\times3\times224\times224$ elementos y que existan nueve
salidas en el orden de \texttt{labels.json}. También exige un segundo tensor de
características para el detector OOD. La pantalla de captura acepta cámara o
galería, ejecuta preprocesamiento e inferencia, navega al resultado y guarda el
escaneo. La pantalla siguiente muestra la fotografía, diagnóstico, confianza,
segunda opción en casos de confianza baja y recomendaciones. Si la imagen se
rechaza, no fuerza una enfermedad: pide acercarse a una hoja con mejor luz.

\subsection{Comprobación ejecutable}

Se instaló el lockfile del repositorio móvil en un clon temporal. La instalación
normal detectó un conflicto entre React Native 0.86.0 y
\texttt{@react-native/jest-preset} 0.86.2; se completó con
\texttt{npm ci --legacy-peer-deps}. Después, \texttt{npm run typecheck} terminó
sin errores y \texttt{npm test -- --runInBand} aprobó 40 suites y 276 pruebas.
Las advertencias de React sobre \texttt{act(...)} no cambiaron el código de
salida. El entorno disponible usa Node 18.19.1 aunque dependencias recientes
solicitan Node 20.19.4 o superior; reproducir la app requiere alinear esa
versión. La instalación también reportó 24 vulnerabilidades de dependencias
(17 moderadas y 7 altas). No se ejecutó una actualización forzada porque
habría cambiado el lockfile del prototipo; deben revisarse antes de distribuir
otra compilación.

La Figura~\ref{fig:prototype-capture} es una captura generada por ejecución, no
una maqueta. El comando carga los pesos públicos del o de los backbones
seleccionados, aplica la cabeza final y escribe un JSON con Top-3, confianza y
latencia CPU. Se usa una imagen del holdout solo después de terminar la
evaluación final.

\begin{figure}[H]
  \centering
  \includegraphics[width=0.94\textwidth]{prototype_capture.png}
  \caption{Captura de la demo de Etapa 2: entrada real y salida producida por el
  pipeline final. La latencia corresponde a CPU de escritorio y no se presenta
  como benchmark Android.}
  \label{fig:prototype-capture}
\end{figure}

La demo se reproduce desde la raíz del clasificador con:

\begin{verbatim}
PYTHONPATH=.:/tmp/doctor-maiz-pydeps python \
  scripts/etapa_2/stage2_experiments.py \
  --output-dir outputs/etapa_2 demo --image /ruta/hoja.jpg
\end{verbatim}

\subsection{Qué modelo está desplegado}

La aplicación móvil conserva EfficientNet-Lite0 de la línea de entrenamiento
principal, no la cabeza congelada evaluada en esta etapa. Hacer el reemplazo
habría requerido exportar conjuntamente backbone, escalado y regresión
one-vs-rest, validar paridad y volver a medir OOD en Android. Ninguna de esas
pruebas se ejecutó, así que el informe no afirma un despliegue que no ocurrió.
El prototipo demuestra que la arquitectura de integración existe y que el
artefacto histórico funciona dentro de ella; la demo demuestra que la selección
experimental produce resultados reproducibles.

Durante una primera consulta, el endpoint de salud respondió
\texttt{\{"status":"ok"\}} y el endpoint de versión informó Android 1.0.0,
código 6. En la verificación final del mismo día, el dominio dejó de resolver
por DNS y la URL de APK sin autenticación devolvió HTTP 404. Por eso el
prototipo se califica como funcional e integrado, pero no como un despliegue
público con disponibilidad comprobada de extremo a extremo. Esta diferencia
reduce la puntuación del criterio y queda como limitación operativa, no como
pendiente oculto.
